% =============================================================================
% Lloyds Commercial Banking Intelligence, group report skeleton
% Built 8 August 2026 from a full read-only crawl of all eight branches.
% Companion document: drafts/project-chronology-2026-08-08.md
%
% HOW TO USE THIS FILE
%   Every heading below is final enough to write into. Under each one there is a
%   \note{...} line saying what goes there, and a % comment giving the owner,
%   the page budget and the artefacts to pull from. Delete each \note as you
%   replace it with prose. Nothing else needs restructuring.
%
% PAGE BUDGET, main matter only (30 pages)
%   1 Introduction ............  3   (10%)
%   2 Background .............   6   (20%)
%   3 Execution ..............   9   (30%)
%   4 Critical Evaluation ....   9   (30%)
%   5 Conclusion .............   3   (10%)
%   Front matter and appendices do not count.
%
% OWNERSHIP
%   Viktor  ch2.3, 2.4, 3.2, 3.4, 3.5, 4.1, 4.3
%   Sneha   1, 2.1, 2.2, 3.1, 3.3, 3.8, 4.6 
%   Vishal  3.7, 4.5
%   Samuel  2.5, 3.6, 4.2, 4.4
% =============================================================================

\documentclass[    author={Team 5},
               supervisor={Prof. Fernando Moreno-Pino},
                   degree={MSc},
                    title={Mining the Public Record for Commercial Banking Intelligence},
                 subtitle={Subtitle},
                     type={},
                     year={2026}]{dissertation}

% One helper, so the skeleton notes are visibly not prose and can be stripped in
% one pass at the end (\renewcommand{\note}[1]{} kills them all).
\newcommand{\note}[1]{\par\noindent\textit{\small [#1]}\par\vspace{2mm}}

\begin{document}

\maketitle
\frontmatter
\makedecl

\tableofcontents
\listoffigures
\listoftables
% \listofalgorithms       % we have no numbered algorithms; leave commented
% \lstlistoflistings      % re-enable only if we float real code listings

% -----------------------------------------------------------------------------
\chapter*{Abstract}
% Compulsory, at most 1 page. Owner: TBC, written LAST.
\note{One page. The problem (Lloyds needs to find BB/SME prospects and risks
among companies that are not already clients), what we built (a 33-month
company-month panel over the Companies House register, four self-labelled
targets, eleven recorded model runs, an explainable ranked shortlist), and what
we found (the structured record predicts; narrative media does not reach
companies of this size at all, and the statutory record reaches them only after
the event). Close with the factual achievement list, in the template's bullet
form: rows of panel built, features engineered, hours of API harvest, runs
recorded, lift achieved. State the outcome view explicitly, including that two
of our headline results are negative and measured.}

% -----------------------------------------------------------------------------
% \chapter*{Summary of Changes}
% Conditional, resubmissions only. Not applicable. Leave commented out.

% -----------------------------------------------------------------------------
\chapter*{Supporting Technologies}
% Compulsory, at most 1 page. Owner: Viktor, with one bullet each from the others.
\note{Bullet list only, grouped as below. Keep it about the project, not about
writing it up. Be explicit about which parts are third-party and where they are
used, and name the things we deliberately did NOT use.}

\begin{itemize}
\item \textbf{Data sources (all public, all free).} Companies House bulk company
      data product and its REST API (companies, officers, filing history,
      charges); Contracts Finder and Find a Tender OCDS feeds; The Gazette;
      HM Land Registry Commercial and Corporate Ownership Data; the Intellectual
      Property Office trade mark register; GLEIF and Wikidata for legal entity
      identifiers; the Guardian Open Platform, NewsAPI, GNews and GDELT for news;
      Adzuna for job adverts; \texttt{yfinance} for listed-company financials.
\item \textbf{Data engineering.} Python, \emph{pandas}, \emph{NumPy},
      \emph{PyArrow} and Parquet for the panel, \emph{DuckDB} for the SQL over
      it and for the entity spine.
\item \textbf{Modelling.} \emph{scikit-learn} (logistic regression, the MLP,
      calibration and metrics), \emph{LightGBM} (gradient-boosted trees),
      \emph{SHAP} for per-company attribution, \emph{rapidfuzz} for name
      resolution.
\item \textbf{Reporting and delivery.} \emph{Matplotlib} for every figure, a
      \texttt{src/models/report.py} module that writes figures and tables
      directly as PDF and \texttt{.tex} for \verb|\includegraphics| and
      \verb|\input|, and a dependency-free static HTML and vanilla-JavaScript
      dashboard.
\item \textbf{Engineering practice.} Git and GitHub with one branch per
      workstream; a run registry that refuses to overwrite a recorded tag; a
      feature hash that makes two runs provably comparable.
\item \textbf{Deliberately not used}, each for a stated reason given in the
      chapter indicated: Streamlit or any served dashboard (\S\ref{sec:delivery});
      fuzzy name matching without a location check (\S\ref{sec:namematch});
      any commercial or paid data source at all.
\end{itemize}

% -----------------------------------------------------------------------------
\chapter*{Notation and Acronyms}
% Optional, 1-2 pages. Owner: TBC. The report uses
% a lot of domain vocabulary and several terms we defined ourselves.
\note{Two tables. The first is ordinary acronyms. The second is the vocabulary
we defined ourselves, which matters more, because a reader who does not know
what we mean by an ``origin'' or a ``gate'' cannot follow Chapter 4 at all.}

\noindent\textbf{Acronyms.}
\begin{quote}
\noindent
\begin{tabular}{lcl}
BB      &:& Business Banking, our smallest client tier (under GBP 3m turnover) \\
SME     &:& Small and Medium Enterprise (GBP 3--25m turnover) \\
LBG     &:& Lloyds Banking Group \\
RM      &:& Relationship Manager, the end user of a shortlist \\
CH      &:& Companies House \\
SIC     &:& Standard Industrial Classification \\
OCDS    &:& Open Contracting Data Standard \\
LEI     &:& Legal Entity Identifier \\
GLEIF   &:& Global Legal Entity Identifier Foundation \\
CCOD    &:& Commercial and Corporate Ownership Data (HM Land Registry) \\
IPO     &:& Intellectual Property Office \\
SHAP    &:& SHapley Additive exPlanations \\
MLP     &:& Multi-Layer Perceptron \\
ROC-AUC &:& area under the receiver operating characteristic curve \\
PR-AUC  &:& area under the precision--recall curve \\
P@$N$   &:& precision at $N$, the hit rate in the top $N$ ranked companies \\
CI      &:& confidence interval \\
\end{tabular}
\end{quote}

\noindent\textbf{Terms we use in a specific sense.}
\begin{quote}
\noindent
\begin{tabular}{lcp{0.62\textwidth}}
snapshot &:& one monthly extract of the Companies House bulk register \\
panel    &:& the company-month table built from 33 snapshots \\
origin   &:& the month $t$ a prediction is made from; features come from $t$, labels from $t{+}1$ onwards \\
horizon $H$ &:& how far forward a label looks, 3, 6 or 12 months depending on target \\
embargo  &:& the gap between the last training origin and the first test origin, so no label window overlaps training \\
as-of join &:& a join that admits a record only if it was \emph{visible} at the origin month, not merely dated before it \\
gate     &:& the visibility rule in an as-of join, written created-lag / satisfied-lag in days \\
base rate &:& the unconditional proportion of positives for a target in a month \\
lift     &:& P@$N$ divided by the base rate, i.e.\ how much better than choosing at random \\
run, tag &:& one recorded execution of the pipeline under a named configuration \\
\end{tabular}
\end{quote}

% -----------------------------------------------------------------------------
\chapter*{Acknowledgements}
% Optional, at most 1 page.
\note{Supervisor, the Lloyds contacts who gave us the brief, and the maintainers
of the public data we relied on entirely. One paragraph is enough.}

% =============================================================================
\mainmatter

% -----------------------------------------------------------------------------
\chapter{Introduction}
\label{chap:introduction}
% 3 pages. Owner: Sneha.

\section{The problem}
\note{A commercial bank can only sell to companies it can see. Lloyds sees its
own book in full and everyone else's hardly at all. The brief asks whether the
public record can be mined to find, among roughly 5.7M registered UK companies,
the ones about to need lending, the ones heading for trouble, and the ones about
to grow, \emph{including companies that do not bank with Lloyds}. State plainly
that this is a search problem over a very large, very sparse population, not a
classification problem over a customer base.}

\section{Why it matters, and to whom}
\note{Three named audiences from the client brief: the Growth team (``not
ours''), the Attrition team (``ours and leaving''), the Maintenance team
(``ours and staying''). Say here that those three are defined by
\emph{counterparty}, not by need, which is why counterparty identity became a
workstream of its own and why the routing in \S\ref{sec:routing} is the piece
that turns a ranked list into something a person can act on.}

\section{Central challenges}
\note{Four that we surmounted and one that we did not, and the fifth is the
honest one. (i) Scale: 5.7M companies x 33 months is 49.6M rows on a laptop.
(ii) Time: the register is a photograph of now, so history had to be
reconstructed from 33 separate snapshots. (iii) Labels: nobody hands you a
label, so every target is self-labelled from a future state of the same
register, which makes leakage the central methodological risk. (iv) Identity:
public sources name companies four different ways and only some of them use the
company number. (v) \textbf{The signal we expected was not there.} Point
forward to N1, N2 and N5. Sneha's draft currently reads as though all four
challenges were surmounted; the fifth needs adding, because it is the finding
the report is built on.}

\section{Aims, objectives and achievements}
\note{Bullet list, per the template, and this is the part examiners read first.
Objective 4 in the current draft says ``full unstructured integration is future
work''. Rewrite it as \emph{achieved with a negative result}, citing N1 and N5.
Measuring that a data source does not work is a completed objective, not a
deferred one, and framing it as deferred gives away the strongest thing we did.
The achievements list should be factual and countable: 33 snapshots, a
49.6M-row panel, 41 engineered features across four families, four
self-labelled targets, 11 recorded model runs across 3 model families, a
25.4-hour API harvest of 558,693 charges, a 284,230-notice Gazette crawl, a
leakage defect found and corrected in four passes, and a shortlist running at
about 150x lift on lending.}

% -----------------------------------------------------------------------------
\chapter{Background}
\label{chap:background}
% 6 pages. Owner: Sneha, Viktor, Sam

\section{Papers: Predicting corporate distress and corporate demand}
The classical approach to predicting company failure is built on financial ratios. The best known example is Altman's Z-score \cite{altman_1968}. It is a statistical model that takes a small set of accounting ratios and combines them using weights estimated by discriminant analysis, into a single number the 'Z-score'. That number places a company in a 'safe', 'grey' or 'distress' zone, a low score flags a firm as likely to fail. Because it is simple, interpretable and only needs figures from published accounts, it remains the reference point for corporate distress models. The difficulty is that it assumes those accounts are available and comparable. Small and medium-sized enterprises (SMEs), which make up the bulk of the UK company register are entitled to file abbreviated or 'micro entity' accounts, which only require a balance sheet and a few short notes, not a full profit and loss statement or detailed financial breakdown \cite{companieshouse_accounts_2026}. For these firms the ratios that drive a Z-score are either missing or too coarse to be useful.

a second line of work responds to this gap by turning to non-financial and behavioural information. Altman, Sabato and Wilson \cite{altman_sabato_wilson_2010} show that non-financial signals materially improve SME default prediction over accounts only models and later credit scoring research finds that behavioral and transactional data can perform as well as or even better than traditional financial variables \cite{berg_etal_2020}. This motivates our own emphasis on register derived behvaiour such as charges, director changes, filing activity rather than on reported financial data alone.

Almost all of this literature predicts distress which firms will fail. The complementary question which firms are about to need finance and ready to take it on is less studied, even though it is the question a commercial lender most wants answered. This 'demand side' framing is what distinguishes our project from conventional default model.

\section{The main data source : UK public company record
}
The backbone of project is the UK public company record maintained by Companies House, which we access in two ways \cite{companieshouse_api}. The first is the bulk product, a snapshot of every company on the register (5,698,276 records at the time of download) conatining static, ready to use fields. The second is the public REST API, which returns richer per company  detail such as officers, filing history, charges but is rate limited. Enriching the full population this way would take more time, so API is applied selectively

The Standard Industrial Classification (SIC) code gives each company an industrial label, which we use for sectorfiltering and grouping \cite{companieshouse_sic}. The accounts category records the type of accounts a firm files and itself is informative, given how many SMEs file abbreviated accounts \cite{companieshouse_accounts_2026}. The mortgage charge counts proxy for secured borrowing and company status records whether a company is active, in liquidation, in adminstration and so on. A practical limitation is that company status is current state, not dated event,os it tells us a firm is insolvent but not when it became so. This becomes important later when we build the labels.

\section{The modelling background: Gradient boosting, ranking, and explanation}
\note{Why gradient-boosted trees on tabular data with heavy missingness; why we
rank rather than classify (a relationship manager works a list, not a decision
boundary), which makes ROC-AUC and precision@$N$ the right pair of metrics and
accuracy the wrong one; and SHAP, both what it computes and why we need it, since
a shortlist without a reason attached to each row cannot be acted on. Keep the
maths light but present: an examiner should be able to read the SHAP figures.}

\section{TBC Leakage and evaluation under time}
\label{sec:leakage-bg}
\note{This section is what makes Chapter 4 legible, and it is worth the space.
Define look-ahead leakage; define the difference between when a fact
\emph{happened} and when it \emph{became visible}, which is the distinction our
whole leakage story turns on; define out-of-time validation and the embargo. If
the reader understands this section, \S\ref{sec:leak} lands as a contribution
rather than as an apology.}

% -----------------------------------------------------------------------------
\chapter{Execution}
\label{chap:execution}
% 9 pages, the largest chapter. Owner: Sneha, Viktor, Sam, Vishal

\section{How the project organised itself and reorganised}
\label{sec:orgs}
Across the projects's lifetime, the team worked in three different organisational setups, each change happened because of what the data showed, not because of preference. The record of who worked on what is visible in the feature branches in repository (Table \ref{tab:branches}.

The first split mirrored the project brief, the brief listed three needs namely growth, attrition and deepening existing relationships, so the team split into three strands to match and early work began along those lines. This arrangement broke within two weeks. The reason was structural, the three needs are three datasets but three questions asked of one dataset, the shared Companies House \cite{companieshouse_api}. Splitting the team by question meant repeating the same cleaning and fetaure work three times, with no one owning the shared base. The strands were merged back together.

The second split was a by data type, one strand handling structured data from Companies House and related public records, and another explored unstructured text such as news articles for company-level signal \cite{guardian_api, newsapi}. This division held for several weeks, because the two halves genuinely needed different tools. It was found that the unstructured strand measured its own coverage and found that, for the population of interest especially SMEs has almost no coverage. Very few companies had any mathcing article or news  mention to draw on. 

That measurement drove the third and final arrangement, which split the team by task rather than by data type. One pair moved into modelling, building the feature matrix, the prediction targets, the classifiers and SHAP explanantions \cite{lundberg_lee_2017} on top of the structured core. The other pair kept working on the unstructured side, news sentiment and the harmonization needed to attach it to establish how much signal, if any, could be recovered there. This is the shape the project held to the end.

Both reorganisations were driven by measurement, not opinion. The first ended when the cost of duplicated work became clear, the second ended when the coverage count for unstructured data came back near zero.

\section{Data acquisition and the base dataset}

\note{What Companies House actually is, what a company is obliged to file and
when, and the single most important property for us: \textbf{the register gives
you status, not the date trouble began}. A company is Active until it is not.
That one fact is what forced every target to be defined as an event inside a
forward window rather than as a state, and it should be stated here so}

\note{The bulk product, 5,698,276 registered companies. The SIC scrape (731
codes) and the mapping of 325 of them onto three target sectors. The Active
filter and the size tiering off \texttt{Accounts.AccountCategory}, with the
regulatory definition of each account category quoted. Result: 1,372,321
companies. State the join key once and for all here: \texttt{CompanyNumber},
zero-padded to eight digits. Mention, in one sentence, that this work was
briefly done twice in parallel and consolidated (Appendix~\ref{appx:chronology}).}
\note{The pivot from cross-section to panel. Acquiring 33 monthly bulk zips
(October 2023 to July 2026) against a verified manifest. The two-pass design,
universe union then per-month extract, and why a naive per-month build mistakes
sector re-coding for dissolution. 49,556,152 rows over about 2.04M companies.
The verification that matters: the panel reproduces the single-vintage build
exactly at June 2026, same row count and same sector counts.
\textbf{Keep the June 2025 footnote}: that month is genuinely absent from the
Companies House server, and rather than interpolating we carry an honest NULL
row through a dense calendar spine. It dictated calendar-aware windows
everywhere downstream and it comes back as a real limitation in
\S\ref{sec:limits}.}

\section{Structured Feature engineering: Companies House and additional source}
The team explored a wide range of public signals and not all of them were kept in the final model, this section records both what we kept and what we discarded along with the reasons. The retained features form four families, all keyed on eight digit \texttt{CompanyNumber}.  Two rules run through every family, every feature is calculated strictly as at the month it belongs to. A feature for month t can only use information that existed in the month t, never anything from the future, which is what  lets \S\ref{sec:targets} define learnable targets. Every feature's definition is stored in a catalogue and turned into a short hash, a finger print of exact input and logic used. This means any two runs can be proven to have used identical outputs, giving reproducibility and preventing silent changes.

\paragraph{Group 1 - The company as it is now}
The first group describes each company as the register currently shows it and it is pulled straight from the Companies Houses bulk file \cite{companieshouse_api} with no extra lookups needed. It covers how old the company is, how many loans it has secured against its assets and how many are still outstanding, a simple debt ratio built from those, whether its accounts are overdue or out of date, a size band which is read from type of accounts it files \cite{companieshouse_accounts_2026}, the industry it works in from its SIC code \cite{companieshouse_sic} and whether it has changed its name. Many of these were prototyped in attrition work, company health read from its status, a dormant flag, a rough size band and late filings warning and checked against early static exploration before comvinign in a single panel.

\paragraph{Group 2 - how the company is changing}
A single snapchat cannot show the changes over the years, and the changes is exactly the signal we mostly want. A company whose secured borrowing has jumped in six months was behaving very differntly from one that has sat still for years, even if the two look identical in the present. So thirty three month history was assembled and measured the change over the last three, six and twelve months. The change in charge counts, in outstanding charges, in the debt ratio and how many separate new-borrowing events happened in the past year was recorded.

One considerable gap in the source shaped the entire design, Companies House never published a snapshot for June 2025 . This is  a permanent gap in their published data, the server was re-checked and the file has never existed, not a download that could be simply repeated. So 'going back three rows' in history to measure the change, could have been landed on May 2025 instead of June 2025. It would have quietly reported a four month change as if it were three, without a warning. Instead every company was laid against a full run of calender months, so the missing month appears as an explicit blank and every look-back is exactly the right number of calendar months, reaching across the gap returns a blank rather than a wrong number. This group also carries the earliest warning sign in the public record such as features that track how logn since a company last filed its accounts or confirmation statement and how long its accounts have been stale. These indicators stay flat while a company is still filling normally and only start to rise once the company stops filing, which usually happens before its status officially changes to insolvency.

\paragraph{Group 3 - winning public work}
This register shows how a company files and borrows, but not whether its winning business. From the Contracts Finder and Find a Tender \cite{contractsfinder, findatender, ocds}, each company's public-sector contract wins were aggregated including the count, their total and average value, and whether it has just won its very first, matched using the Companies House number carried in contract record. Part-way through the window the rules changed. When Procurement Act 2023 came into force on 24 February 2025, new public contracts stopped being published on Contracts Finder and began appearing on Find a Tender \cite{findatender}, while Contracts Finder kept only older notices. So information from both were combined to avoid false drop in wins that has nothing to do with companies themselves. Every win is counted only up to the month in question, so 2026 award never lands on 2024 row.

\paragraph{Group 4 - who the company banks with}
For companies that have secured loans, lender names attached to each charge was read using Charges API \cite{companieshouse_api} , every name was sorted using a clear ordered rule list, into either "a Lloyds brand" (Lloyds Bank, Bank of Scotland, Halifax, HBOS and related names \cite{lbg_brands} or "a rival". This turns a raw public record into necessary information of who the company banks with, the very axis the client's growth and attrition questions are asked on. Two more relationship signals from Officiers and Filling History records were added, recent director appointments and resignations, and new-share (SH01) filings as a sign that company is raising money.

\subsection{Signals we tried but left out}
Several signals were built and measured, then deliberately left out. The financial data route failed on coverage, the market data feeds only carry companies listed on stock exchange and almost none of the private small businesses were listed, so the data simple did not describe the population. The news route failed the same way for a different reason, when how many of the targeted companies had any matching news was measured, the answer was almost none, so there was no signal to add. Both were honest dead ends found by measurement. Furthermore, the unstructured data for companies based on signals were added as an enrichment layer.



\note{The four families and why each exists. Level features from the register
as it stands. Delta features at 3, 6 and 12 months, computed on a dense calendar
spine rather than a positional lag, so the June 2025 hole yields NULLs instead of
silently wrong four-month deltas. Filing-behaviour features closing the ``stopped
filing'' gap. Public contract features. Finish with the feature catalogue and
the \texttt{feature\_hash}, which is what lets us prove two runs used identical
inputs. Full catalogue in Appendix~\ref{appx:features}.}

\S\ref{sec:targets} does not have to argue for it from scratch.
\begin{itemize}
    \item notebooks/13\_panel\_deltas.ipynb Static features (all companies). Computable directly from the bulk file for every company: company
age, secured charge counts and a derived debt ratio, accounts-overdue and accounts-stale flags, size tier,
and previous-name / name-change indicators.
\item notebooks/14b\_lender\_charges.ipynb Relationship (lender) signals. For companies with outstanding charges, I query the Companies
House Charges API and read the persons entitled lender names, classifying each as Lloyds-linked
(Lloyds Bank, Bank of Scotland, Halifax, HBOS, and related brands) or competitor. This turns a public
charge record into a coarse banking-relationship signal.
\item notebooks/14\_contracts\_asof.ipynb Contract-activity signals. From the Contracts Finder OCDS feed I aggregate public-sector award
wins to the company level (contracts won, total and average value), joined on the GB-COH Companies
House number.
\item notebooks/06\_director\_changes.ipynb Director and growth-filing signals. From the Officers and Filing History APIs I derive recent director appointments and resignations, and recent share-allotment (SH01) filings as a capital-raise proxy.
\item attrition\_analysis.ipynb
\end{itemize}

\subsection{Public procurement, joined as of each month}
\note{Both OCDS bulk feeds, Contracts Finder and Find a Tender, unioned across
the Procurement Act 2023 publishing cutover of 24 February 2025. Explain why:
using either feed alone produces a decay artefact that is purely regulatory and
would have been read as a real collapse in contract activity. 2,633 cross-source
duplicate awards found and removed. 173k awards over 44.6k companies, and note
the scale gap against the earlier API-paged version (2,154 awards over 1,317
companies), which is a concrete argument for bulk feeds over per-company API
calls.}

\subsection{Name resolution, and why it is dangerous}
\label{sec:namematch}
\note{The suppliers with no company number. Normalised name plus postcode,
measured at 88\% precision on name alone and 92\% with postcode over about 100k
labelled records. Two design decisions worth defending: match against the
\emph{full} register rather than our narrow three-sector universe, because
matching against the narrow set manufactures false positives; and keep the
extended set as a separate A/B directory so the strict set stays the default for
anything a relationship manager sees. Cross-reference Samuel's false-positive
examples in \S\ref{sec:sources-eval}, since they are the cautionary evidence for
this whole subsection.}

\subsection{Counterparty identity: the Charges harvest}
\note{Why lender identity is the axis the client brief is actually defined on.
The harvest: 158,684 companies, 558,693 charges, 25.4 hours, zero 404s, staged
and resumable. The ordered regex taxonomy collapsing free-text lender names onto
institutions, classifying 77.5\% of charge-lender rows, and the honest handling
of the remaining 22.5\% as a visible ``unclassified'' population rather than a
silent one. The thirteen lender features. \textbf{The point-in-time gate belongs
here as a design decision} (a charge is admitted only once it was visible, not
once it was created); its failure and repair are the story of \S\ref{sec:leak}.
Also record the \texttt{switching} label: built, 797 positives against a
pre-registered 1,000-positive stop rule, therefore not trained, and shipped as
three rule-based feeds instead. Pre-registering a threshold and then honouring
it when it went against us is worth its own sentence.}

\section{Targets and the leakage-safe split}
\label{sec:targets}
\note{The section that makes everything after it credible. Features from origin
$t$, labels strictly from $t{+}1$ to $t{+}H$. The four targets with horizons and
base rates: Lending Readiness (3m, 0.25--0.31\%), Credit Risk Exposure (6m,
0.30--0.41\%), Voluntary Exit (6m, 6.7--8.5\%), Growth Signal (12m,
2.1--2.2\%). The out-of-time split, the embargo, and \texttt{FIRST\_FULL\_ORIGIN}.
Negative downsampling at 10:1 with the keep rate carried per row so
recalibration still works. \textbf{Include both mid-build corrections}, they are
the most instructive part: distress redefined from ``non-Active'' (90\% of which
was benign strike-off) to genuine insolvency only, and voluntary exit redefined
from a state at $t{+}H$ to an event anywhere in the window, roughly doubling the
base rate, because about half of strike-off proposals are withdrawn but still
mark a real distress moment.}

\section{The modelling pipeline as software}
\note{Frame this as reproducibility engineering, not as plumbing.
\texttt{RunConfig}, the three-family model registry, the recalibration step, and
a \texttt{record\_run} that refuses to overwrite an existing tag. Staged
execution via a script rather than a notebook, and the honest reason: the
machine kept rebooting under long jobs. The \texttt{src/} module map, one line
each. Point at Appendix~\ref{appx:runs} for the run registry.}

\subsection{Baseline}

\section{Unstructured data}
\note{Samuel's survey. One short paragraph per source: Contracts Finder and
Find a Tender, The Gazette, Land Registry CCOD, the IPO trade mark register,
GLEIF and the LEI system, Adzuna, and the news APIs. \textbf{For each one give a
single extra sentence saying whether it keys on the company number or on the
company name.} That one distinction predicts, almost perfectly, which sources
worked for us and which did not, and setting it up here lets Chapter 4 pay it
off in a line. Include \texttt{yfinance} and listed-company data here as a
documented dead end (26 hand-picked listed firms; 0.00\% of our universe holds a
ticker) rather than giving it a section of its own.}

\subsection{The unstructured attempt I: narrative media}
\note{viktor-vishal/unstructured-data-lab. Method only; the result is
\S\ref{sec:unstructured}. The 96-company balanced stress-test grid, four size
tiers by three sectors by eight, explicitly non-representative so the rare cells
have a measurable $n$. The API bake-off and why GDELT was dropped early. The
six levels of query construction from bare name to domain restriction. The
disambiguation rule: exact phrase plus corroboration on town, sector or context,
with town mandatory for single common-word names. The positive control design,
and why a project reporting a zero needs one.}

\subsection{The unstructured attempt II: the statutory record}
\note{Method only. The Gazette has no company-name search, only category, date
and postcode, so notices are pulled in bulk and matched back through the company
number printed in each one. One request per ten seconds, cached per page and
resumable, about 7.6 hours. 284,230 notices extracted, 274,904 after dropping
ceremonial noise. The five-stage distress ladder and the 52 company-level
features. \textbf{Set this section up as the deliberate contrast with the
previous one}: same goal, same team, opposite coverage profile, and the reason
is that one source keys on a name and the other keys on a number.}

\subsection{Spine crosswalk}
\note{This section would cover samuel/spine-crosswalk. we never used this but if we have space we can talk about it}

\section{Dashboard}
\label{sec:delivery}
\note{The dashboard, Vishal and Samuel. Precompute-then-serve: one Python build
pass flattens the widened universe and the Gazette features into a single store,
and a static page with no framework reads it. Defend the two decisions that look
like shortcuts and are not: no server, so it survives a demo on someone else's
laptop; and the store is a \texttt{.js} assignment rather than \texttt{.json}
because browsers block \texttt{fetch()} on \texttt{file://} URLs. Seven declared
sources of which one is live, with the other six shown as ``not connected yet''
rather than as zero, which is an honesty decision worth naming. Screenshots in
Appendix~\ref{appx:dashboard}.}

\subsection{Routing the shortlist by counterparty}
\label{sec:routing}
\note{The piece that connects the model to the brief. The five-way partition of
the universe by lender relationship (current LBG client 11,003; lapsed 13,421;
borrowing elsewhere and never ours 64,531; charge held but lender unclassified
14,653; no charge ever 1,305,676), and the mapping of the first three onto the
client's Maintenance, Attrition and Growth teams. Keep the unclassified
population visibly separate rather than folding it in: it is the taxonomy gap
made honest at company level, and hiding it is the one place the lender
dictionary's limits could become a client-facing error. Each row carries its top
three SHAP reasons.}

\section{Engineering practice}
\note{Feature definitions live in a shared package rather than being re-implemented per notebook, so the training
matrix matches the exploratory notebooks exactly. Version control uses a branch per workstream, and
the modelling run records its configuration and metrics so that a re-run is a diff rather than a mystery}

% -----------------------------------------------------------------------------
\chapter{Critical Evaluation}
\label{chap:evaluation}
% 9 pages, the largest chapter. Owner: Sneha, Viktor, Sam, Vishal

\section{Modelling design eval + results}

\subsection{Evaluation design}
\note{State the rules before the numbers, so no result later looks chosen after
the fact. Why ranking rather than classification. Why ROC-AUC and precision@$N$
together, and what each one actually answers: ROC-AUC is a global ranking score
that survives the negative downsampling, precision@$N$ is what a relationship
manager experiences. Why per-origin rather than pooled, which \S\ref{sec:errors}
then justifies with a measurement. What the bootstrap intervals resample.}

\subsection{Model comparison across eleven runs}
\note{The main results table: 11 runs, 4 targets, 3 model families, every cell
with an interval. \texttt{\input} the generated table. Read out the headline
comparisons and, importantly, the non-comparisons: where two runs are
indistinguishable, say so, rather than reporting the larger number.}

\subsection{Is the sample variable we are predicting balanced? Recall comparison}

\subsection{A leakage defect, found and fixed}
\label{sec:leak}
\note{The strongest methodological section in the report, and per the 8 August
decision it stays a section here rather than becoming its own chapter. Tell it
in four acts. (i) The lender run returns precision@100 = 1.000 on lending
against a 0.271\% base rate, while ROC-AUC barely moves; explain why that
combination is a leakage signature rather than a win, because a genuinely better
feature set lifts the global ranking and the top of the list together. (ii) The
cause: the as-of gate admitted a charge on \texttt{created\_on}, but 94.7\% of
charges are delivered to Companies House after they are created, so the model
could see a charge the register had not yet recorded, and the register is what
defines the label. (iii) The obvious fix was not enough: gating on
\texttt{delivered\_on} removed about half of it. (iv) Calibrating properly showed
that 18 of the 21 days of the interim margin were the \emph{satisfaction} clock,
not a visibility lag, so one gate had been widened to pay for a different gate
being wrong. Include the retraction: the extract-drift hypothesis was
satisfying, and it was worth 0.11 days of 21, and it ran the opposite way to
prediction. Also state that the original verification test inherited the bug
rather than catching it, because it replayed the harvest against the same wrong
clock, and that a single-month version of the rewritten test would still have
passed. Figures: the gate progression on lending P@500, the calibration ladder,
the lag-error plots.}

\subsection{Does counterparty identity help? (N4)}
\note{The verdict, read off the most conservative gate against the deterministic
control, and the reasoning that makes it more than a null. The apparent gain
scales monotonically with how much margin the gate allows, in every model family;
one model getting greedier is a story about that model, three families moving in
lockstep is a story about the data. The mechanism sentence the section turns on:
the lending label is computed from the very register the gate governs, so
loosening visibility leaks the label directly, whereas insolvency comes from a
separate register and shows no gradient at all under the same loosening. Two
targets, opposite behaviour, one mechanism, and we did not design that natural
experiment. Then the answer with its caveats: at the tightest gate tested, the
lender features give a boosted model nothing and give the linear model a small
but real gain, which locates the finding in \emph{encoding} rather than in
information, because the base feature set already carries mortgage-charge counts.
Both caveats travel with it: the tightest gate tested is not a gate proved
clean, and the most impressive precision number the project produced is one we
are throwing away on purpose.}

\subsection{Sneha doc sections}

\subsection{What our error bars are worth}
\label{sec:errors}
\note{Present this as a result, not a footnote, because it is one. Two findings.
(i) We validated our confidence intervals against a second, independent source of
variation, refitting under different seeds, and they held for LightGBM and
logistic and failed for the MLP, where a seed term dominates. In one case two
non-overlapping intervals were reassuring us about a difference whose sign flips
across seeds. State the resulting rule: no MLP difference in this project below
about 0.02 ROC-AUC is interpretable. (ii) The paired bootstrap: comparing two
runs by separate intervals is the wrong test when both are evaluated on the same
companies, and resampling once and differencing tightens the lending interval
about five-fold and moves it from inconclusive to clearly excluding zero.
(iii) Pooled precision@$N$ is not an average of its months; it lands outside the
range of the two months it pools in 9 of 12 rows, in both directions, and the
worst row lands 0.108 below the worse of its own months. So per-origin figures
are the honest ones, and it is also what a relationship manager receives, one
month at a time.}

\section{The unstructured result: narrative against statutory (N1, N5)}
\label{sec:unstructured}
\note{Vishal and Samuel's headline, and after \S\ref{sec:leak} the sharpest
section in the report. It needs one voice, not two. The coverage table across
NewsAPI, Guardian, GNews and GDELT; the raw hit rate of 2.8--6.9\%; the
verification step; and then 0 of 96 companies and 0 of 83 articles surviving
disambiguation, with the Tesco control returning 1,209 to prove the plumbing
works. Name the actual collisions, they are memorable and they make the point
better than the number does: ``Baroness'' returning Michelle Mone, ``Ncc''
returning the acronym rather than the company, ``Peco'' returning a Pennsylvania
utility. \textbf{Include the superseded measurement, attributed}: an earlier run
(NB08, Vishal) reported 95.8\% coverage, which was a broadcast bug collapsing
per-company results to a constant per segment-by-sector cell, and NB09 by the
same author rebuilt the query and superseded it. Say who did which, since the
individual reports have to be consistent with this one. Then the turn: the Gazette covers 106,664 companies, and only 1,033 of
them are inside the Active universe, because a company in liquidation is on its
way out of Active status. Also report, as a method decision rather than a gap,
the 43\% of notices carrying no company number that were discarded rather than
name-matched. Close on the synthesis, which is the thesis of the chapter:
journalism does not reach small companies at all; the statutory record reaches
them completely but only once something has already gone wrong, which is
precisely when a bank no longer needs a prediction.}

\section{What the model is worth commercially}
\note{The shortlists, from the control model with the growth fix, deliberately
not from a lender run. Per held-out month, with lift: lending 0.43 and 0.41 at
about 150x, insolvency 0.16 and 0.14 at about 45x, growth 0.22 at about 10x,
voluntary exit 0.85 and 0.78 at about 10x. \textbf{Instruct the reader to read
the lift column, not the precision column}, and use voluntary exit as the worked
example of why: it looks like the best list and is the weakest, because roughly
one company in twelve exits anyway. State that the four lists are almost
disjoint, 398 unique companies across 400 rows, and that scores are comparable
within a target and month but not across targets. Show one company end to end
with its score and its top three SHAP reasons, because that is the actual
deliverable.}

\section{What each public source was worth (N2, N3)}
\label{sec:sources-eval}
\note{Samuel's evaluation, and the payoff for the keying sentence planted in
\S2.5. One table: source, what it keys on, coverage of the spine, verdict.
Land Registry 5.25\% and the widest SME reach we found; trade marks 1.04\%;
LEI 0.10\%; ticker 0.00\%; Adzuna one advert matched in a thousand.
\textbf{Report the LEI correction openly and attributed}: Samuel first measured
9.14\%, traced it to an orphan join counting companies that were never in the
spine, and corrected it to 0.10\%. Finding an inflated number in your own work
and correcting it is worth more than the number was, and attributing it is what
makes the individual reports line up with this one. Same for the name-matching false positives, an Oxford ``Amazon
Ltd'' credited with AWS contracts and an ``Ace Tax Ltd'' credited with taxi
work, which are the concrete justification for the postcode requirement in
\S\ref{sec:namematch}. Finish with N3: 7.2\% of firms lost a bank in the 24
months before distress against a 5.5\% baseline, so losing a bank describes
attrition and does not predict it.}

\section{Competitor market intelligence}
\note{\textbf{TBC, decide by 16 August.} The Charges harvest is a map of UK SME
secured lending, and this section reads it without a model, a label or a
precision caveat. Four analyses: the league table by sector and segment (NatWest
16.4\%, LBG 12.1\%, Barclays 12.0\%, HSBC 11.9\%, so the big four hold about half
the register and NatWest holds a third more of it than we do); stock against
flow, to show which lenders are losing ground; the switching-flow matrix between
banks; and wallet-share concentration on the LBG book. Carry the limits with it:
secured lending only, 7.8\% of the universe, 22.5\% of lender names
unclassified, and a charge is a proxy for a banking relationship rather than the
relationship itself.}

\section{Limitations}
\label{sec:limits}
\note{Collect them honestly and do not soften them; this section makes the
chapter stronger. The June 2025 gap in the source and its real consequence for
any one-month lag at the July 2025 origin. The growth defect: nine features that
were 100\% NULL in every training origin and populated in test, so the
41-feature count was a fiction for that target and it is now a 27-feature model
that says what it is. Voluntary exit precision@$N$ is partly arbitrary because
982 of the top 1000 scores are exact ties, so nobody should quote it to three
digits. MLP seed instability, per \S\ref{sec:errors}. The two-company residual
that the rewritten verification still fails on the historical lender panel, and
why it nearly went unnoticed. Four run manifests that record the wrong gate
value. One misleadingly named feature that stays misleadingly named because
renaming it would invalidate the comparability of every recorded run. Secured
lending is only a proxy for a banking relationship. And the population question:
Active-only caps the Gazette at 1,033 companies, so a statutory-informed model
would need a wider population than the one we modelled.}

% -----------------------------------------------------------------------------
\chapter{Conclusion}
\label{chap:conclusion}
% 3 pages. Owner: TBC.

\section{Contributions}
\note{Re-summarise, in the order they matter. A reproducible 33-month panel over
the whole UK register with a run registry behind it. Four self-labelled targets
and a ranked, explainable shortlist that routes to the three teams the client
named. Five measured negatives that most projects would have left unmeasured.
And a methodological contribution that is genuinely transferable beyond this
domain: the distinction between when a fact happened and when it became visible,
and what it costs to get that wrong.}

\section{Status against the aims}
\note{Objective by objective, with the chapter that evidences each. Say plainly
what is working and what is not: the structured pipeline is complete and
verified; the unstructured integration exists as a delivery layer with one of
seven sources live and a measured reason for the other six; the model scores
are joined into the dashboard; and the competitor market intelligence is either
done or explicitly out, per the decision taken on 16 August.}

\section{Further work}
\note{Focus on the interesting openings rather than on a wish list. The gate
gradient has not flattened at the tightest margin we tested, so the leak-free
intercept may be lower still and a tighter rung would settle it. Persist
held-out prediction vectors per model family, since the verdict rests partly on
logistic gains we cannot paired-test. Widen the modelling population to include
Fading and Insolvent companies so the statutory signal can inform the model
rather than decorate it. And the enrichment loop that we parked: rank the top
500 per target, enrich all of them, re-run, now that the ranking half costs one
command and only the API cost per company is unknown.}

% =============================================================================
\backmatter

\bibliography{bibtex.bib}
% Owner: Sneha, Viktor, Sam, Vishal

% -----------------------------------------------------------------------------
\appendix
% Appendices do not count toward the 30 pages, so anything that is evidence
% rather than argument belongs here. Examiners are not obliged to read them,
% which means nothing load-bearing may live here and nothing here may be the
% only statement of a result.

\chapter{Project chronology by branch}
\label{appx:chronology}
\note{The full dated record, from \texttt{drafts/project-chronology-2026-08-08.md}:
eight branches, every notebook, what it produced, and why it was kept,
superseded or stopped. Include the abandoned work, the duplicated filtering, the
attrition branch, the discredited 95.8\% news figure and the listed-company
notebook that was never executed. This is also where the individual
contributions are made auditable.}

\chapter{Sector selection and SIC mapping}
\label{appx:sic}
\note{The 325 SIC codes mapped onto the three target sectors, the size tiering
from account category with each regulatory definition, and the resulting
segment-by-sector counts.}

\chapter{The feature catalogue}
\label{appx:features}
\note{All 41 control features and the 13 lender extensions: name, family,
definition, source, and whether it is level, delta, contract or lender.
\texttt{\input} the generated catalogue rather than retyping it.}

\chapter{Target definitions in full}
\label{appx:targets}
\note{The exact SQL-level definition of each of the four targets plus the
untrained switching label, the horizons, the origin lists for train and test,
and the base rate per origin.}

\chapter{Run registry and reproducibility}
\label{appx:runs}
\note{Every recorded run with its tag, configuration, feature hash, git SHA and
full metric table with intervals; and the commands to reproduce any of them.
Include the four manifests that record the wrong gate value, with the true
values, so an auditor reading manifests alone is not misled.}

\chapter{The dashboard}
\label{appx:dashboard}
\note{Screenshots of the landing page, the watchlist, a company page with the
model indices filled, and the shortlist explorer with the counterparty routing.
One short paragraph on how to run it: open the file.}

% =============================================================================
\end{document}